\section{Evidence from a Published Leaderboard}\label{sec:external}

Everything so far concerns a percentile constructed here, on a panel assembled here, so the effect may belong to that construction rather than to leaderboards in general. HELM Lite publishes its full accuracy table at every versioned release, archives the earlier releases, and reports a pool-relative headline \citep{liang2022helm}. The same evidence can therefore be read at fourteen points in time while the set of models being compared grows underneath it.

If a leaderboard moves, the first thing to rule out is a revised suite. Across releases v1.0.0 through v1.13.0 the ten scenario columns are identical, with nothing added, dropped, or renamed apart from a relabelled column header at v1.10.0. The same 24 models appear in all fourteen releases, and their 240 published scenario scores are bit-identical across the entire sequence. Over the same span the evaluated pool grows from thirty models to ninety-one.

Against that frozen evidence, all twenty-four of the constant models see their published Mean win rate change, by 0.1232 on average in absolute value and by as much as 0.2178 for Mixtral 8x7B, which falls from 0.7276 to 0.5098. Most of that movement is arithmetic: Mean win rate is the fraction of head-to-head comparisons a model wins, so a fixed model faces a larger and later field as the pool grows, and the change correlates with initial standing at $-0.766$.

The movement that matters is elsewhere. If the drift were a monotone transform of the frozen evidence it would be harmless, since a monotone rescaling preserves every order. It is not. Of the 24, twenty-two fall and two rise, and five of the 276 pairs among them reverse between the first release and the last. \autoref{fig:helm} plots each model's rank among the constant set at the first release against the last, so a reversed pair is a crossing: the five are countable on the page, and each occurs while the evidence beneath it is fixed. \citet{roelofs2019meta} asks whether reusing a test set degrades the evidence and finds little of it across a hundred competitions; here the evidence is fixed and the reference set is what moves. The published numbers are correct as defined at every release. What changed was the pool, and the published order followed the pool. \citet{singh2025leaderboard} report a comparable pattern from a different mechanism, deprecation policy rather than a coverage boundary changing which models remain available, and in both the reference set rather than the evidence moves the ranking.

The control comes from the same organisation. HELM Capabilities runs on the same infrastructure and the same release cadence, and its pool grows comparably, from twenty-two models to sixty-eight over sixteen releases. Its headline is a Mean score, an absolute average over scenarios rather than a fraction of comparisons won. The same 22 models appear in every release, their 110 published scenario scores are bit-identical, and none of the twenty-two headline values move at all. No pair reorders, at the endpoints or ever. That prediction, zero drift under an absolute headline with the same pool growth, was stated before the test: the mechanism is the pool-relative aggregation, not the evaluation practice, the suite, or the organisation. \autoref{fig:helm} shows the nine models whose order changed, named.

\begin{figure}[H]\centering
% generated by src.figures_tikz -- do not edit by hand
\begin{tikzpicture}
  \begin{axis}[causalre body, xmin=-0.5, ymin=0, ymax=1, ytick={0,0.25,0.5,0.75,1}, xlabel={release}, xmax=13.5,
    ylabel={published headline value}, title={(a)}]
    \addplot[cMUTE, smooth, line width=0.5pt, forget plot] coordinates {(0,0.0552) (1,0.0677) (2,0.0800) (3,0.1000) (4,0.0891) (5,0.0833) (6,0.0765) (7,0.0775) (8,0.0744) (9,0.0731) (10,0.0695) (11,0.0679) (12,0.0663) (13,0.0645)};
    \addplot[cMUTE, smooth, line width=0.5pt, forget plot] coordinates {(0,0.6207) (1,0.6290) (2,0.5925) (3,0.5766) (4,0.5509) (5,0.5317) (6,0.5124) (7,0.5034) (8,0.4951) (9,0.4851) (10,0.4675) (11,0.4599) (12,0.4539) (13,0.4392)};
    \addplot[cMUTE, smooth, line width=0.5pt, forget plot] coordinates {(0,0.9621) (1,0.9645) (2,0.9575) (3,0.9574) (4,0.9345) (5,0.9183) (6,0.9220) (7,0.9154) (8,0.9148) (9,0.9089) (10,0.8853) (11,0.8809) (12,0.8778) (13,0.8671)};
    \addplot[cMUTE, smooth, line width=0.5pt, forget plot] coordinates {(0,0.8345) (1,0.8419) (2,0.8250) (3,0.8213) (4,0.7909) (5,0.7733) (6,0.7683) (7,0.7584) (8,0.7560) (9,0.7454) (10,0.7212) (11,0.7171) (12,0.7109) (13,0.6982)};
    \addplot[cMUTE, smooth, line width=0.5pt, forget plot] coordinates {(0,0.2517) (1,0.2613) (2,0.2450) (3,0.2532) (4,0.2364) (5,0.2217) (6,0.2075) (7,0.2030) (8,0.1961) (9,0.1924) (10,0.1843) (11,0.1811) (12,0.1780) (13,0.1723)};
    \addplot[cMUTE, smooth, line width=0.5pt, forget plot] coordinates {(0,0.3483) (1,0.3516) (2,0.3250) (3,0.3319) (4,0.3109) (5,0.2933) (6,0.2796) (7,0.2734) (8,0.2650) (9,0.2598) (10,0.2483) (11,0.2436) (12,0.2414) (13,0.2329)};
    \addplot[cMUTE, smooth, line width=0.5pt, forget plot] coordinates {(0,0.6586) (1,0.6645) (2,0.6250) (3,0.6085) (4,0.5891) (5,0.5683) (6,0.5460) (7,0.5370) (8,0.5261) (9,0.5215) (10,0.5058) (11,0.4997) (12,0.4951) (13,0.4819)};
    \addplot[cMUTE, smooth, line width=0.5pt, forget plot] coordinates {(0,0.2172) (1,0.2258) (2,0.2125) (3,0.2340) (4,0.2091) (5,0.1983) (6,0.1840) (7,0.1804) (8,0.1744) (9,0.1700) (10,0.1629) (11,0.1602) (12,0.1576) (13,0.1517)};
    \addplot[cMUTE, smooth, line width=0.5pt, forget plot] coordinates {(0,0.0345) (1,0.0452) (2,0.0500) (3,0.0723) (4,0.0636) (5,0.0600) (6,0.0529) (7,0.0521) (8,0.0500) (9,0.0474) (10,0.0451) (11,0.0440) (12,0.0430) (13,0.0411)};
    \addplot[cMUTE, smooth, line width=0.5pt, forget plot] coordinates {(0,0.0862) (1,0.0968) (2,0.1000) (3,0.1191) (4,0.1073) (5,0.1017) (6,0.0927) (7,0.0930) (8,0.0906) (9,0.0904) (10,0.0860) (11,0.0840) (12,0.0820) (13,0.0784)};
    \addplot[cMUTE, smooth, line width=0.5pt, forget plot] coordinates {(0,0.2034) (1,0.2161) (2,0.2050) (3,0.2128) (4,0.1982) (5,0.1850) (6,0.1751) (7,0.1719) (8,0.1663) (9,0.1635) (10,0.1555) (11,0.1530) (12,0.1506) (13,0.1450)};
    \addplot[cMUTE, smooth, line width=0.5pt, forget plot] coordinates {(0,0.4379) (1,0.4484) (2,0.4175) (3,0.4149) (4,0.3873) (5,0.3667) (6,0.3444) (7,0.3383) (8,0.3286) (9,0.3252) (10,0.3130) (11,0.3067) (12,0.3031) (13,0.2918)};
    \addplot[cMUTE, smooth, line width=0.5pt, forget plot] coordinates {(0,0.8207) (1,0.8323) (2,0.8250) (3,0.8255) (4,0.7909) (5,0.7683) (6,0.7615) (7,0.7491) (8,0.7403) (9,0.7234) (10,0.7040) (11,0.6979) (12,0.6921) (13,0.6792)};
    \addplot[cMUTE, smooth, line width=0.5pt, forget plot] coordinates {(0,0.7724) (1,0.7806) (2,0.7375) (3,0.7234) (4,0.6909) (5,0.6633) (6,0.6462) (7,0.6344) (8,0.6262) (9,0.6171) (10,0.5986) (11,0.5903) (12,0.5847) (13,0.5698)};
    \addplot[cMUTE, smooth, line width=0.5pt, forget plot] coordinates {(0,0.3759) (1,0.3871) (2,0.3525) (3,0.3511) (4,0.3218) (5,0.3083) (6,0.2914) (7,0.2890) (8,0.2826) (9,0.2810) (10,0.2721) (11,0.2680) (12,0.2641) (13,0.2535)};
    \addplot[cSF, smooth, line width=0.5pt, forget plot] coordinates {(0,0.3241) (1,0.3258) (2,0.3025) (3,0.3064) (4,0.2818) (5,0.2667) (6,0.2546) (7,0.2495) (8,0.2421) (9,0.2399) (10,0.2306) (11,0.2263) (12,0.2246) (13,0.2168)};
    \addplot[cSF, smooth, line width=0.5pt, forget plot] coordinates {(0,0.5138) (1,0.5226) (2,0.4750) (3,0.4681) (4,0.4418) (5,0.4217) (6,0.4003) (7,0.3918) (8,0.3827) (9,0.3746) (10,0.3600) (11,0.3538) (12,0.3490) (13,0.3357)};
    \addplot[cSF, smooth, line width=0.5pt, forget plot] coordinates {(0,0.5069) (1,0.5226) (2,0.4975) (3,0.5021) (4,0.4745) (5,0.4517) (6,0.4283) (7,0.4201) (8,0.4084) (9,0.4003) (10,0.3844) (11,0.3788) (12,0.3723) (13,0.3580)};
    \addplot[cSF, smooth, line width=0.5pt, forget plot] coordinates {(0,0.3310) (1,0.3387) (2,0.3125) (3,0.3170) (4,0.2909) (5,0.2767) (6,0.2605) (7,0.2537) (8,0.2461) (9,0.2399) (10,0.2294) (11,0.2251) (12,0.2222) (13,0.2145)};
    \addplot[cSF, smooth, line width=0.5pt, forget plot] coordinates {(0,0.5034) (1,0.5097) (2,0.4700) (3,0.4660) (4,0.4345) (5,0.4183) (6,0.3974) (7,0.3904) (8,0.3827) (9,0.3785) (10,0.3649) (11,0.3609) (12,0.3560) (13,0.3446)};
    \addplot[cSF, smooth, line width=0.5pt, forget plot] coordinates {(0,0.7276) (1,0.7387) (2,0.6950) (3,0.6787) (4,0.6527) (5,0.6217) (6,0.5947) (7,0.5823) (8,0.5708) (9,0.5626) (10,0.5413) (11,0.5320) (12,0.5254) (13,0.5098)};
    \addplot[cSF, smooth, line width=0.5pt, forget plot] coordinates {(0,0.6862) (1,0.6935) (2,0.6725) (3,0.6553) (4,0.6364) (5,0.6133) (6,0.5932) (7,0.5836) (8,0.5748) (9,0.5677) (10,0.5522) (11,0.5450) (12,0.5405) (13,0.5264)};
    \addplot[cSF, smooth, line width=0.5pt, forget plot] coordinates {(0,0.7759) (1,0.7839) (2,0.7700) (3,0.7809) (4,0.7491) (5,0.7283) (6,0.7189) (7,0.7026) (8,0.6930) (9,0.6811) (10,0.6637) (11,0.6598) (12,0.6549) (13,0.6436)};
    \addplot[cSF, smooth, line width=0.5pt, forget plot] coordinates {(0,0.7828) (1,0.7935) (2,0.7575) (3,0.7362) (4,0.7164) (5,0.6950) (6,0.6715) (7,0.6586) (8,0.6494) (9,0.6398) (10,0.6208) (11,0.6119) (12,0.6058) (13,0.5889)};
  \end{axis}
  \begin{axis}[causalre body, xmin=-0.5, ymin=0, ymax=1, ytick={0,0.25,0.5,0.75,1}, xlabel={release}, xmax=15.5,
    xshift=0.5\textwidth, yticklabels={,,}, title={(b)}]
    \addplot[cNYC, smooth, line width=0.5pt, forget plot] coordinates {(0,0.5513) (1,0.5513) (2,0.5513) (3,0.5513) (4,0.5513) (5,0.5513) (6,0.5513) (7,0.5513) (8,0.5513) (9,0.5513) (10,0.5513) (11,0.5513) (12,0.5513) (13,0.5513) (14,0.5513) (15,0.5513)};
    \addplot[cNYC, smooth, line width=0.5pt, forget plot] coordinates {(0,0.5222) (1,0.5222) (2,0.5222) (3,0.5222) (4,0.5222) (5,0.5222) (6,0.5222) (7,0.5222) (8,0.5222) (9,0.5222) (10,0.5222) (11,0.5222) (12,0.5222) (13,0.5222) (14,0.5222) (15,0.5222)};
    \addplot[cNYC, smooth, line width=0.5pt, forget plot] coordinates {(0,0.5907) (1,0.5907) (2,0.5907) (3,0.5907) (4,0.5907) (5,0.5907) (6,0.5907) (7,0.5907) (8,0.5907) (9,0.5907) (10,0.5907) (11,0.5907) (12,0.5907) (13,0.5907) (14,0.5907) (15,0.5907)};
    \addplot[cNYC, smooth, line width=0.5pt, forget plot] coordinates {(0,0.5487) (1,0.5487) (2,0.5487) (3,0.5487) (4,0.5487) (5,0.5487) (6,0.5487) (7,0.5487) (8,0.5487) (9,0.5487) (10,0.5487) (11,0.5487) (12,0.5487) (13,0.5487) (14,0.5487) (15,0.5487)};
    \addplot[cNYC, smooth, line width=0.5pt, forget plot] coordinates {(0,0.6531) (1,0.6531) (2,0.6531) (3,0.6531) (4,0.6531) (5,0.6531) (6,0.6531) (7,0.6531) (8,0.6531) (9,0.6531) (10,0.6531) (11,0.6531) (12,0.6531) (13,0.6531) (14,0.6531) (15,0.6531)};
    \addplot[cNYC, smooth, line width=0.5pt, forget plot] coordinates {(0,0.6741) (1,0.6741) (2,0.6741) (3,0.6741) (4,0.6741) (5,0.6741) (6,0.6741) (7,0.6741) (8,0.6741) (9,0.6741) (10,0.6741) (11,0.6741) (12,0.6741) (13,0.6741) (14,0.6741) (15,0.6741)};
    \addplot[cNYC, smooth, line width=0.5pt, forget plot] coordinates {(0,0.6653) (1,0.6653) (2,0.6653) (3,0.6653) (4,0.6653) (5,0.6653) (6,0.6653) (7,0.6653) (8,0.6653) (9,0.6653) (10,0.6653) (11,0.6653) (12,0.6653) (13,0.6653) (14,0.6653) (15,0.6653)};
    \addplot[cNYC, smooth, line width=0.5pt, forget plot] coordinates {(0,0.6343) (1,0.6343) (2,0.6343) (3,0.6343) (4,0.6343) (5,0.6343) (6,0.6343) (7,0.6343) (8,0.6343) (9,0.6343) (10,0.6343) (11,0.6343) (12,0.6343) (13,0.6343) (14,0.6343) (15,0.6343)};
    \addplot[cNYC, smooth, line width=0.5pt, forget plot] coordinates {(0,0.5647) (1,0.5647) (2,0.5647) (3,0.5647) (4,0.5647) (5,0.5647) (6,0.5647) (7,0.5647) (8,0.5647) (9,0.5647) (10,0.5647) (11,0.5647) (12,0.5647) (13,0.5647) (14,0.5647) (15,0.5647)};
    \addplot[cNYC, smooth, line width=0.5pt, forget plot] coordinates {(0,0.6086) (1,0.6086) (2,0.6086) (3,0.6086) (4,0.6086) (5,0.6086) (6,0.6086) (7,0.6086) (8,0.6086) (9,0.6086) (10,0.6086) (11,0.6086) (12,0.6086) (13,0.6086) (14,0.6086) (15,0.6086)};
    \addplot[cNYC, smooth, line width=0.5pt, forget plot] coordinates {(0,0.6569) (1,0.6569) (2,0.6569) (3,0.6569) (4,0.6569) (5,0.6569) (6,0.6569) (7,0.6569) (8,0.6569) (9,0.6569) (10,0.6569) (11,0.6569) (12,0.6569) (13,0.6569) (14,0.6569) (15,0.6569)};
    \addplot[cNYC, smooth, line width=0.5pt, forget plot] coordinates {(0,0.6786) (1,0.6786) (2,0.6786) (3,0.6786) (4,0.6786) (5,0.6786) (6,0.6786) (7,0.6786) (8,0.6786) (9,0.6786) (10,0.6786) (11,0.6786) (12,0.6786) (13,0.6786) (14,0.6786) (15,0.6786)};
    \addplot[cNYC, smooth, line width=0.5pt, forget plot] coordinates {(0,0.6416) (1,0.6416) (2,0.6416) (3,0.6416) (4,0.6416) (5,0.6416) (6,0.6416) (7,0.6416) (8,0.6416) (9,0.6416) (10,0.6416) (11,0.6416) (12,0.6416) (13,0.6416) (14,0.6416) (15,0.6416)};
    \addplot[cNYC, smooth, line width=0.5pt, forget plot] coordinates {(0,0.6177) (1,0.6177) (2,0.6177) (3,0.6177) (4,0.6177) (5,0.6177) (6,0.6177) (7,0.6177) (8,0.6177) (9,0.6177) (10,0.6177) (11,0.6177) (12,0.6177) (13,0.6177) (14,0.6177) (15,0.6177)};
    \addplot[cNYC, smooth, line width=0.5pt, forget plot] coordinates {(0,0.5737) (1,0.5737) (2,0.5737) (3,0.5737) (4,0.5737) (5,0.5737) (6,0.5737) (7,0.5737) (8,0.5737) (9,0.5737) (10,0.5737) (11,0.5737) (12,0.5737) (13,0.5737) (14,0.5737) (15,0.5737)};
    \addplot[cNYC, smooth, line width=0.5pt, forget plot] coordinates {(0,0.4437) (1,0.4437) (2,0.4437) (3,0.4437) (4,0.4437) (5,0.4437) (6,0.4437) (7,0.4437) (8,0.4437) (9,0.4437) (10,0.4437) (11,0.4437) (12,0.4437) (13,0.4437) (14,0.4437) (15,0.4437)};
    \addplot[cNYC, smooth, line width=0.5pt, forget plot] coordinates {(0,0.3759) (1,0.3759) (2,0.3759) (3,0.3759) (4,0.3759) (5,0.3759) (6,0.3759) (7,0.3759) (8,0.3759) (9,0.3759) (10,0.3759) (11,0.3759) (12,0.3759) (13,0.3759) (14,0.3759) (15,0.3759)};
    \addplot[cNYC, smooth, line width=0.5pt, forget plot] coordinates {(0,0.4785) (1,0.4785) (2,0.4785) (3,0.4785) (4,0.4785) (5,0.4785) (6,0.4785) (7,0.4785) (8,0.4785) (9,0.4785) (10,0.4785) (11,0.4785) (12,0.4785) (13,0.4785) (14,0.4785) (15,0.4785)};
    \addplot[cNYC, smooth, line width=0.5pt, forget plot] coordinates {(0,0.3966) (1,0.3966) (2,0.3966) (3,0.3966) (4,0.3966) (5,0.3966) (6,0.3966) (7,0.3966) (8,0.3966) (9,0.3966) (10,0.3966) (11,0.3966) (12,0.3966) (13,0.3966) (14,0.3966) (15,0.3966)};
    \addplot[cNYC, smooth, line width=0.5pt, forget plot] coordinates {(0,0.6091) (1,0.6091) (2,0.6091) (3,0.6091) (4,0.6091) (5,0.6091) (6,0.6091) (7,0.6091) (8,0.6091) (9,0.6091) (10,0.6091) (11,0.6091) (12,0.6091) (13,0.6091) (14,0.6091) (15,0.6091)};
    \addplot[cNYC, smooth, line width=0.5pt, forget plot] coordinates {(0,0.5990) (1,0.5990) (2,0.5990) (3,0.5990) (4,0.5990) (5,0.5990) (6,0.5990) (7,0.5990) (8,0.5990) (9,0.5990) (10,0.5990) (11,0.5990) (12,0.5990) (13,0.5990) (14,0.5990) (15,0.5990)};
    \addplot[cNYC, smooth, line width=0.5pt, forget plot] coordinates {(0,0.5291) (1,0.5291) (2,0.5291) (3,0.5291) (4,0.5291) (5,0.5291) (6,0.5291) (7,0.5291) (8,0.5291) (9,0.5291) (10,0.5291) (11,0.5291) (12,0.5291) (13,0.5291) (14,0.5291) (15,0.5291)};
  \end{axis}
\end{tikzpicture}

\caption{The nine HELM Lite models whose published order changed between the first release
and the last, ranked among the 24 whose 240 scenario scores are bit-identical throughout;
a crossing is a pair whose order reversed, one hue per group of mutually reversing models,
and there are five. The control: across sixteen HELM Capabilities releases behind an
absolute headline, with the pool growing from 22 to 68, none of the 22 constant models'
headline values moved and no pair reversed.}
\label{fig:helm}
\end{figure}

This need not be read as an error on HELM's part: the absolute headline on its own Capabilities leaderboard is the point rather than an embarrassment. A pool-relative headline is a joint statement about a model and about the population evaluated beside it. Two such numbers need not be comparable even when every measurement behind them is shared. Read later as statements about the models alone, they move standing between providers for reasons belonging to neither. Restoring comparability across moments is the classical equating problem \citep{kolen2014}, which an absolute scale estimated jointly across benchmarks avoids entirely \citep{epoch2026rosetta}. The panel and the leaderboard show the same effect.
